\documentclass[10pt,a4paper]{amsart}
\usepackage[margin=19mm]{geometry}
\usepackage{amsmath,amssymb,mathtools}
\usepackage[scr=rsfs,cal=euler]{mathalfa}
\usepackage{xfrac}
\usepackage[unicode,hidelinks,bookmarks=false]{hyperref}
\newcommand{\ie}{i.e.\ }
\newcommand{\cf}{cf.\ }
\newcommand{\resp}{respectively\ }
\newcommand{\Subsection}{\subsection}
\providecommand{\nobreakdash}{\nobreak\hbox{-}}
\setlength{\emergencystretch}{2em}
\multlinegap0pt
\usepackage{mathtools}
\usepackage{amssymb}
\usepackage{enumerate}
\usepackage{mdframed}
\usepackage{xcolor}
\newtheorem{corollary}{Corollary}
\newtheorem{lemma}{Lemma}
\usepackage{cleveref}
\crefname{corollary}{Corollary}{Corollarys}
\crefname{lemma}{Lemma}{Lemmas}
\newcommand{\C}{\mathbf C}
\newcommand{\Cc}{\mathcal C}
\newcommand{\brak}{[\![\hbar]\!]}
\newcommand{\cbrak}{(\!(\hbar)\!)}
\newcommand{\mr}{\mathrm}
\title{\vspace{-1cm}Differential graded categories in holomorphic symplectic geometry}
\author{Borislav Mladenov}
\date{Source: arXiv:2604.06630v1 (CC BY 4.0)}
\begin{document}
\maketitle

\noindent\emph{Source and licence.} \vspace{-1cm}Differential graded categories in holomorphic symplectic geometry. Version arXiv:2604.06630v1 (CC BY 4.0), distributed by arXiv under the Creative Commons Attribution 4.0 International licence, http://creativecommons.org/licenses/by/4.0/, as recorded in the arXiv licence metadata for this exact version and shown on the abstract page https://arxiv.org/abs/2604.06630v1. Full licence notice: \texttt{licenses/arxiv-2604.06630-CC-BY-4.0.txt}.\par\smallskip

\noindent\textit{Editorial scope.} Counted unit: the article's lemma lemma (source0.tex lines 1235--1237). The article's complete original proof of it is delivered with it (source0.tex lines 1238--1252, one \texttt{proof} environment of 15 lines). No mathematics is written, repaired or re-derived: the statement and the proof are the article's own lines, in the article's own notation, and no retained line is substituted or re-flowed.  Retained uncounted as the source-local context the proof needs: lines 636--638; lines 1217--1219; lines 1227--1230.  Nothing was omitted from any retained span, so the package carries no omission record.  No retained line contains a citation, so the wrapper carries no bibliography and no bibliography entry.  No retained line contains a figure environment, an image-inclusion command or drawing code, so no image is delivered and the package declares no asset dependency.  Editorial formatting only, in addition: an AMS \texttt{amsart} preamble that restates the article's own declarations and macros used by the retained lines, namely the environments \texttt{corollary}, \texttt{lemma} on separate counters, together with the standard packages those lines need.  The retained environments print in this document as Corollary 1, Lemma 1 in order of appearance, whereas the article prints the same items with the numbering of its own full document.  The complete native source of the article as distributed in the arXiv e-print for this exact version is bundled unchanged as \texttt{sources/198031/source0.tex}.
\begin{corollary}\label{a2dg}
Let $\Cc$ be a flat, proper finite dg category over an integral domain $\mr{R}$ with generic point $\eta$. Assume the $\mr{R}$-module $\mr{HH}^2_c(\mr{H}\Cc,m_2)$ is torsion-free, where $m_2$ is the composition in $\Cc$. Then $\Cc$ is formal iff $\Cc_\eta$ is formal.
\end{corollary}

\begin{corollary}\label{lags}
Let $\mr{X}/\C$ be a holomorphic symplectic manifold. Let $\mathfrak{L}$ be a Solomon-Verbitsky collection. Then the holomorphic Solomon-Verbitsky category $\mathcal{DQ}_\mathfrak{L}$ is formal.
\end{corollary}

 \begin{lemma}\label{freelemma}
 Let $\iota : \C \xhookrightarrow{} \C\brak$ be the inclusion of the central fibre and let $$\mr{C} \in \mr{Perf}\big(\mr{Spec}\big(\C\brak\big)\big).$$ Suppose that for all $i\in \mathbf{Z}$ we have $$\mr{dim}_\C\big(\mr{H}^{i}(\iota^{*}\mr{C})\big) = \mr{dim}_{\C\cbrak}\big(\mr{H}^{i}(\C\cbrak\otimes_{\C\brak} \mr{C})\big).$$
 Then the cohomology $\mr{H}(\mr{C})$ is free (of finite rank) over $\C\brak$.
 \end{lemma}

\begin{lemma}
Let $\mathfrak{L}$ be a Solomon-Verbitsky collection and let $\widetilde{\mathcal{D}}\subset \widetilde{\mathcal{D}}_\mathfrak{L}$ be a finite full dg subcategory of $\widetilde{\mathcal{DQ}}_\mathfrak{L}$. Then it is formal and quasi-isomorphic to the trivial formal deformation of $\iota^*\widetilde{\mathcal{D}}$, where $\iota :\C\to \C\brak$ is the inclusion of the central fibre.
\end{lemma}

\begin{proof}
Let $\mathcal{D}=\mr{loc}\left(\widetilde{\mathcal{D}}\right)$ and $\mathcal{D}^\mr{vir}\subset \mathcal{DR}^{\mr{vir}}_\mathfrak{L}$ be the full subcategory spanned by the objects of $\mathcal{D}$. We know $\mathcal{D}$ is formal by \cref{lags}. The degeneration of the local-to-global $\mr{Ext}$ spectral sequences implies that $\mr{H}\widetilde{\mathcal{D}}$ is free of finite type over $\C\brak$. Hence we think of it as a (flat) formal deformation of $\mr{H}\iota^*\widetilde{\mathcal{D}}\cong \iota^*\mr{H}\widetilde{\mathcal{D}}$. Since $\mr{H}\widetilde{\mathcal{D}}$ is finite and its morphism spaces are free of finite rank over $\C\brak$, upper semi-continuity gives an inequality \begin{equation}\label{eq1}\mr{dim}_{\C\cbrak}\left(\mr{HH}^{p,q}\left(\mr{H}\mathcal{D}\right)\right)\le \mr{dim}_{\C}\left(\mr{HH}^{p,q}\left(\iota^*\mr{H}\widetilde{\mathcal{D}}\right)\right).\end{equation} On the other hand, the spectral sequence argument shows that there is a multiplicative filtration $\mr{F}$ on $\iota^*\mr{H}\widetilde{\mathcal{D}}$ such that $$\mr{gr}_\mr{F}\left(\iota^*\mr{H}\widetilde{\mathcal{D}}\right) \cong  \mr{H}\mathcal{D}^\mr{vir}.$$ The completed Rees deformation associated to $\mr{F}$ gives a deformation with central fibre $\mr{H}\mathcal{D}^{\mr{vir}}$. Thus
\begin{equation}\label{eq2}\mr{dim}_{\C\cbrak}\left(\mr{HH}^{p,q}\left(\mr{H}\mathcal{D}\right)\right)=\mr{dim}_{\C}\left(\mr{HH}^{p,q}\left(\mr{H}\mathcal{D}^\mr{vir}\right)\right)\ge \mr{dim}_{\C}\left(\mr{HH}^{p,q}\left(\iota^*\mr{H}\widetilde{\mathcal{D}}\right)\right).\end{equation}
Combining \eqref{eq1} and \eqref{eq2} with \cref{freelemma}, we get the compactly supported (of second kind) Hochschild cohomology $\mr{HH}_c^\bullet\left(\mr{H}\widetilde{\mathcal{D}}\right)$ is free over $\C\brak$. Now formality of $\mathcal{D}$ and \cref{a2dg} imply $\widetilde{\mathcal{D}}$ is formal.\\
Abbreviate the Hochschild cohomology of the graded categories as follows: $\mr{HH} \coloneqq \mr{HH}\big(\iota^*\mr{H}\widetilde{\mathcal{D}}\big)$, $\mr{HH}_{\hbar} \coloneqq \mr{HH}\big(\mr{H}\widetilde{\mathcal{D}}\big)$ and let $\mr{HH}_{(\hbar)}$ be the localisation of $\mr{HH}_\hbar$ at $\hbar$. The reduction $\mr{mod}\,\hbar$ gives a map $\mr{HH}_{\hbar} \to \mr{HH}$ and since $\mr{HH}_\hbar$ is free over $\C\brak$, taking a section of the reduction map, we get $\mr{HH}_\hbar \cong \mr{HH}\brak$.\\
  We have an isomorphism of graded categories $$\varphi_{\hbar} :\mr{H}\widetilde{\mathcal{D}}\otimes_{\C\brak}\C\cbrak = \mr{H}\mathcal{D}\cong \mr{Ind}_{\C\cbrak/\C}\left(\mr{H}\mathcal{D}^{\mr{vir}}\right).$$ Write $m_{\hbar}$ and $m_{\mr{dR}}$ for the compositions of $\mr{H}\widetilde{\mathcal{D}}$ and $\mr{H}\mathcal{D}^{\mr{vir}}$, respectively. Then by definition
  \begin{equation}\label{mult}
    m_{\hbar} = \varphi_{\hbar}^{-1}m_{\mr{dR}}(\varphi_{\hbar},\varphi_{\hbar}).
  \end{equation}
 Differentiating \eqref{mult} with respect to $\hbar$ gives
 \begin{equation}\label{mult2}
 m_{\hbar}' = -\mr{d}_{\hbar}(\varphi_{\hbar}^{-1}\varphi_{\hbar}'),
 \end{equation}
 where $\mr{d}_{\hbar}$ is the Hochschild differential of the graded category $\mr{H}\widetilde{\mathcal{D}}$, i.e. $$[m_\hbar']=0\text{ in }\mr{HH}^{2,0}_{(\hbar)}.$$ Since we have a free $\C\brak$-module, we may write $$m_{\hbar} = m + m_{r}\hbar^{r} + \cdots,\, r \ge 1,$$ hence we see that the left side of \eqref{mult2} is $$rm_{r}\hbar^{r-1}+(r+1)m_{r+1}\hbar^{r} + \cdots.$$ Hence we deduce that $$m_{r}+(r+1)/r\cdot m_{r+1}\hbar + (r+2)/r\cdot m_{r+2}\hbar^{2} +\cdots$$ is an $\hbar$-torsion class in $\mr{HH}^{2,0}_\hbar$, lifting the class $[m_r] \in \mr{HH}^{2,0}$, so it must vanish. Thus $[m_r]=0$ too. Letting $m_r =\mr{d}(\psi^{r})$, we define an automorphism $$\mr{H}\widetilde{\mathcal{D}}\to \mr{H}\widetilde{\mathcal{D}}$$ by acting as identity on objects and $\mr{id}-\psi^{r}\hbar^r$ on morphisms. It kills $m_r$. By induction we get $\psi^{i}$ for all $i \ge r$, the infinite composition  $$\psi_\hbar = \big(\big(\mr{id}-\psi^{r}\hbar^r\big)\circ \big(\mr{id}-\psi^{r+1}\hbar^{r+1}\big)\circ \cdots \big)$$ makes sense and we have $\psi_\hbar^{-1}m_\hbar(\psi_\hbar,\psi_\hbar) = m$, showing the triviality of $m_\hbar$.
\end{proof}

\end{document}
